\chapter*{Lời Tựa}
\addcontentsline{toc}{chapter}{Lời Tựa}
\markboth{Lời Tựa}{Lời Tựa}

\begin{truyen}{Lời Tựa}{Symmetry is a mirror, not a prison.}
\chuhoa{S}{ống thêm một đoạn, người ta bớt tin vào mấy câu chốt hạ cho đẹp. Đời không chạy theo kiểu có là sướng, mất là xong, đúng là hay, sai là bỏ. Nhiều khi được một thứ thì cũng đang mất một thứ. Nhiều khi sai một cú đau điếng lại làm mình tỉnh hơn cả chục lời khuyên.}
\textit{(The older we grow, the more we see that life does not run on simple labels. Having is not always fullness. Losing is not always the end. Being right is not yet complete if it lacks humanity, and being wrong can become the starting point of a wiser way to live.)}

Cuốn sách này viết theo trục \textbf{các cặp đối xứng} của đời sống. Không phải để bắt người đọc chọn phe, mà để nhìn cho ra hai mặt của một chuyện, cái giá của mỗi bên, và cái ngu của việc lao đầu quá đà về một phía.
\textit{(This book is rewritten around the axis of \textbf{symmetrical pairs} in human life. Each pair is not meant to force us to choose a side, but to help us see the balance, cost, and limits within every choice.)}

Người học có thể đọc để soi lại mình. Người dạy học có thể lấy làm chất liệu nói chuyện trong lớp. Người đang mệt, đang bực, đang rối có thể dùng nó như một chỗ dừng để xếp lại đầu óc và lòng dạ.
\textit{(In reading this book, students may find lessons for themselves, teachers may find situations that invite classroom dialogue, and those tired by life may find a pause to rearrange their inner world.)}
\end{truyen}

\begin{baihoc}
Nói gọn: đời không phẳng. Cứ nghĩ chuyện gì cũng một màu là sớm muộn cũng trả giá.
\textit{(Life philosophy does not lie in choosing a side for the world, but in understanding both faces of an issue so we can live with less tension, less shallowness, and less haste.)}
\end{baihoc}

\begin{ghinhoanh}
\textbf{Short English to remember:}

Symmetry teaches balance.

Wisdom begins when labels become questions.
\end{ghinhoanh}

\begin{tuhocnhanh}
1. Trong các cặp đối xứng của cuốn sách, em đang mất cân bằng nhất ở cặp nào? \textit{(Among these symmetrical pairs, where do you feel most out of balance?)}

2. Em thường sống theo một cực đoan nào mà ít khi nhìn lại mặt còn lại? \textit{(Which extreme do you tend to live by without looking at the other side?)}
\end{tuhocnhanh}

\begin{flushright}
\textbf{Nguyễn Văn Sang}\\
\textit{GV THPT Nguyễn Hữu Cảnh - TP HCM}
\end{flushright}

\ngancach
